\subsection{Mô tả chi tiết các trường hợp sử dụng}
\subsubsection{Kiểm soát vào/ra thành viên của trường (mem\_io)}
\begin{adjustbox}{width=1\textwidth,center=\textwidth}
\small
\begin{tabular}{|>{\centering\arraybackslash}m{3.5cm}|m{13.2cm}|}
\hline
\textbf{Use-case ID} & mem\_io \\ \hline
\textbf{Use-case name} & Kiểm soát vào/ra thành viên của trường \\ \hline
\textbf{Use-case overview} & Kiểm soát quá trình vào/ra của các thành viên trường (sinh viên, giảng viên, nhân viên) bằng thẻ nhận diện ID, tích hợp với cơ sở hạ tầng của trường và dịch vụ xác thực trung tâm. \\ \hline
\textbf{Actors} & Các thành viên trường HCMUT, HCMUT\_SSO (Phụ), HCMUT\_DATACORE (Phụ), Hệ thống IoT tại chỗ. \\ \hline
\textbf{Preconditions} & 
1. Hệ thống Quản lý Bãi đỗ xe Thông minh cho Khuôn viên Trường Đại học đang hoạt động và bãi đỗ còn chỗ trống. \newline
2. Thành viên thực hiện thủ tục vào/ra có thẻ nhận diện hợp lệ. \newline
3. Kết nối đến HCMUT\_SSO và HCMUT\_DATACORE phải khả dụng để xác thực theo thời gian thực. \newline
4. Thiết bị phần cứng của barrier (đầu đọc RFID, servor motor, sensors...) đang hoạt động. \\ \hline
\textbf{Trigger} & Một thành viên trường dùng thẻ nhận diện ID chạm vào đầu đọc thẻ tại cổng barrier. \\ \hline
\textbf{Steps} & 
1. Người dùng chạm thẻ nhận diện ID vào đầu đọc RFID. \newline
2. Hệ thống đọc thông tin từ thẻ nhận diện ID và xác thực người dùng thông qua dịch vụ HCMUT\_SSO. \newline
3. Hệ thống truy xuất thông tin người dùng và quyền sử dụng bãi đỗ xe của trường từ HCMUT\_DATACORE. \newline
4. Hệ thống xác thực yêu cầu thông cổng (kiểm tra hạn sử dụng của quyền sử dụng bãi đỗ xe, chống chuỗi nhập xuất bất thường như hai lần vào liên tục,...). \newline
5. Hệ thống ghi record thông cổng (ID người dùng, timestamp, ID cổng, hạn sử dụng quyền sử dụng bãi đỗ,...) trên log. \newline
6. Hệ thống cập nhật (tăng 1) tổng số lần gửi xe vào HCMUT\_DATACORE và giá trị xuất/nhập của thành viên. \newline
7. Hệ thống gửi lệnh đến servor motor controller của barrier để mở cổng. \newline
8. Cảm biến nhận diện chuyển động của phương tiện qua barrier. \newline
9. Hệ thống gửi lệnh đến servor motor controller của barrier để đóng cổng sau khi cảm biến nhận diện phương tiện đã qua cổng. \\ \hline
\end{tabular}
\end{adjustbox}

\begin{adjustbox}{width=1\textwidth,center=\textwidth}
\small
\begin{tabular}{|>{\centering\arraybackslash}m{3.5cm}|m{13.2cm}|}
\hline
\textbf{Post conditions} & Quá trình thông cổng được ghi log lại, trạng thái nhập xuất của thành viên được cập nhật và cổng barrier đã đóng . \\ \hline
\textbf{Exception flow} & 
1. \textbf{Thẻ nhận diện bất hợp lệ/hết hạn:} Hệ thống từ chối thông cổng, hiện lỗi trên màn hình tại cổng, barrier tiếp tục đóng, phát cảnh báo cho nhân viên túc trực. \newline
2. \textbf{Lỗi kết nối:} Nếu không thể kết nối đến HCMUT\_SSO và HCMUT\_DATACORE, hệ thống sẽ chuyển sang sử dụng dữ liệu local và vào chế độ Manual (xử lý bởi spec\_occ). \newline
3. \textbf{Lỗi chuỗi nhập xuất bất thường:} Nếu có thành viên yêu cầu vào cổng nhưng hệ thống đã ghi nhận người này Inside (1 ví dụ trong các trường hợp lỗi này), hệ thống sẽ ghi lại log và phát cảnh báo cho nhân viên túc trực. \newline
4. \textbf{Lỗi barrier:} Nếu cảm biến barrier nhận diện có vật thể (người, phương tiện) cản trợ việc đóng cổng (barrier đang mở) thì barrier sẽ tiếp tục mở. Nếu thời gian mở hơn 15s thì sẽ phát cảnh báo. \\ \hline
\end{tabular}
\end{adjustbox}